\silentchapter{4 Landskapet er dynamisk}{landskapet er dynamisk}
\prop{4}{}{Landskapet er dynamisk.}
\prop{4.1}{a}{Seleksjonstrykka kviler aldri\footnote{Teknologisk gjeld og endra kravspesifikasjonar syter for at trykka fluktuerer konstant. Krav til responstid og datakonsistens omformar domenet utan varsel.}.}
\prop{4.11}{t}{Av 3.1 og 4.1: landskapet er ei uroleg flate\footnote{Nettverkstopologien er aldri statisk; nodar fell ut, og bandbreidda snevrast inn.}.}
\prop{4.111}{t}{Haugar kan stige, søkke, flytte seg, bifurkere eller smelte saman.}
\prop{4.1111}{t}{Ein posisjon som var optimalt tilpassa under eitt sett av vilkår, kan i neste augeblink vere suboptimal\footnote{Ein arkitektur som er optimal for tusen brukarar vert suboptimal for ein million; tilpassinga er lokal i både rom og skala. Det synkrone paradigmet knekk under asynkron last.}.}
\prop{4.11111}{t}{Optimal tilpassing er alltid lokal i tid.}
\prop{4.11112}{o}{Ein ny teknologi flyttar grensa for det forbodne\footnote{Nye abstraksjonslag opnar forbodne regionar av kompleksitet som før krasja grunna minnehandtering. Distribuert konsensus gjer tidlegare umoglege transaksjonar trivielle.}.}
\prop{4.111121}{o}{Ein ny marknad rekonfigurerer kva som vert belønna\footnote{Brukaråtferda si dreining mot mobile flater straffar informasjonstettleik og premierer minimalistisk interaksjon.}.}
\prop{4.111122}{o}{Eit nytt materiale utvidar mengda av moglege operasjonar\footnote{Overgangen frå relasjonelle til grafbaserte datastrukturar mogleggjer nye typologiar av oppslag.}.}
\prop{4.111123}{t}{Kvart av desse er eit inngrep i landskapet.}
\prop{4.11113}{o}{Den rullande medianen av høgde delt på breidde fell systematisk over fem hundreår.}
\prop{4.111131}{o}{Proporsjonen 1.36 er ikkje meir ergonomisk enn 1.88: begge er innanfor menneskeleg sitteområde.}
\prop{4.111132}{o}{Det er ei formdrift, ikkje ei funksjonsdrift.}
\prop{4.2}{o}{Endringa er diskontinuerleg.\footnote{Eldredge \& Gould (1972)\footnote{Eit paradigmeskifte i programmering manifesterer seg nøyaktig slik: stase broten av brå topologiske restruktureringar. Overgangen frå monolitt til mikrotjenester slettar det gamle kartet på eit sekund.}}}
\prop{4.21}{o}{Lange periodar med stase vert avbrotne av brå topologiske skifte\footnote{Heile klassar av designmønster vert forelda over natta når underliggande køyre-miljø vert oppdaterte.}.}
\prop{4.211}{o}{Det typiske forløpet er: eksplosiv radiasjon inn i ein nyopna region, fylgd av gradvis konvergens mot nye attraktorar\footnote{Nye bibliotek utløyser eksplosiv radiasjon av løysingar, før dei best eigna mønstera vert standardiserte. Eit nytt grensesnitt skaper vill flora før konvensjonane set seg.}.}
\prop{4.2111}{t}{Episodisk diskontinuitet er ikkje brot med den underliggande dynamikken.}
\prop{4.21111}{t}{Ho er den dynamikken sin synlege signatur.}
\prop{4.21112}{o}{At endringsraten er diskontinuerleg, ikkje bimodalt fordelt, presiserer dette.}
\prop{4.21113}{o}{Det modernistiske brotet etter 1900 manifesterer seg som eit rekurrensmatrisemønster der ingen periode før 1900 liknar nokon periode etter 1900.}
\prop{4.211131}{o}{Formhistoria gjentar seg ikkje på tvers av dette brotpunktet.}
\prop{4.3}{t}{Landskapet har minne.\footnote{Arthur (1994)\footnote{Den eksisterande databasestrukturen er eit minnespor som avgrensar kva for nye tenester som rasjonelt kan implementerast. Tidlegare avgjerder ligg innbakte i koden som spøkelsestrykk.}}}
\prop{4.31}{o}{Kvar realisert form etterlét eit spor som modifiserer seleksjonstrykka for den neste\footnote{Eit publisert API bind framtida; fjerning av eit endepunkt øydelegg det omkringliggande økosystemet.}.}
\prop{4.311}{t}{Forma og landskapet endrar kvarandre gjensidig.}
\prop{4.3111}{t}{All formgjeving er omformgjeving: agenten startar aldri frå ein tom posisjon\footnote{Utviklaren byrjar aldri med blanke ark; rammeverket er allereie lada med stiavhengige seleksjonstrykk. Gammal kode dikterer rammene for ny kode.}.}
\prop{4.31111}{t}{At agenten startar frå ein posisjon, tyder at ho arvar alt denne posisjonen impliserer.}
\prop{4.311111}{t}{Ho arvar alle naboregionane som er tilgjengelege, og alle dei som er utelukka.}
\prop{4.311112}{t}{Utgangspunktet er aldri nøytralt.}
\prop{4.31112}{t}{Stiavhengigheit er difor eit strukturelt trekk ved formhistoria, ikkje ein anomali\footnote{Valet av eit asynkront paradigme låser systemet inn i eit spor der visse synkrone mønster vert utilgjengelege utan massiv refaktorering. Arven stenger dører.}.}
\prop{4.311121}{o}{At mahogni okkuperte 16 av 16 norske stolar i perioden 1825–1849, etter å ha vore fråverande i den føregåande tilgjengelege perioden, er stiavhengig seleksjonskollaps i si klåraste empiriske form.}
\prop{4.4}{o}{Formrommet ekspanderer kumulativt\footnote{API-økosystem ekspanderer monotont; fjerning av funksjonar skjuler berre kompleksiteten, medan det underliggande konseptrommet veks.}.}
\prop{4.41}{o}{Nye regionar vert busette; gamle vert sjeldan heilt forlatne\footnote{Tidlegare versjonar av ein protokoll må vedlikehaldast for bakoverkompatibilitet; lag på lag med historie eksisterer samstundes.}.}
\prop{4.411}{o}{Det kumulative konvekse hylstervolumet veks monotont over alle tilgjengelege periodar, utan ein einaste tilbakegang.}
\prop{4.4111}{t}{Eit fast funksjonsenvelop ville krevja metning.}
\prop{4.41111}{o}{Nyheitsraten fell ikkje monotont mot metning; han hentar seg inn att når nye materiale kjem inn.}
\prop{4.41112}{t}{Det tilstøytande moglege er eit aktivt ekspanderande rom, ikkje eit tømmande lager.\footnote{Kauffman (1993)\footnote{Å openkjeldekode eit bibliotek er ikkje å tømme eit rom, men å gjere eksponentielt fleire tilstøytande posisjonar tilgjengelege for eksterne agentar.}}}
\prop{4.41113}{t}{Materialstraumar er ein hovudmekanisme for landskapsendring.}
\prop{4.411131}{o}{Kvart nytt materiale opnar nye regionar og stenger andre\footnote{Introduksjonen av distribuert logikk opnar for global skala, men stenger for atomiske, synkrone tilstandsgarantiar.}.}
\prop{4.411132}{o}{Materialstraumen 1500–2025 syner klåre seleksjonsbølgjer.}
\prop{4.411133}{o}{Kvar av dei er ei lokal kohorthending, ikkje gradvise overgangar mellom funksjonelt overlegne alternativ.}
\prop{4.5}{o}{Når eitt seleksjonstrykk vert dominant, kollapsar landskapet mot éin attraktor\footnote{Viss berre prosesseringsfart vektast og minnekostnad ignorerast, kollapsar systemarkitekturen mot ubrukbar rigiditet. Eit absolutt krav om tryggleik frys systemet i statisk uvirksamheit.}.}
\prop{4.51}{t}{Eit einsidig trykk som køyrer over alle andre er ikkje ei forklaring av form.}
\prop{4.511}{t}{Det er eit fråvær av forklaring.}
\prop{4.5111}{t}{Diversiteten i det realiserte formrommet er ein funksjon av mangfaldet i dei aktive trykka\footnote{Eit robust distribusjonsnettverk krev balanse mellom latens, gjennomstrøyming og feiltoleranse; dominans av éin parameter drep systemet.}.}
\prop{4.51111}{t}{Ein monokultur i formrommet er eit symptom på at eitt trykk har vorte uforholdsmessig dominant\footnote{Kloning av brukar-grensesnitt på tvers av bransjar indikerer at frykta for kognitiv friksjon har slått ut alle andre estetiske omsyn.}.}
